\subsection{苏斯林集的分离}

定理 2. 设 X 是可度量化空间。若给定 X 的两两不相交的苏斯林子空间的一个序列 \((\mathbf{X}_{n})\) ，则存在 X 中两两不相交的波莱尔集的一个序列 \((\mathbf{B}_{n})\) ，使得对每个 n 有 \(X_{n} \subset B_{n}\) 。

证明基于两个引理：

引理 4. 设 \((\mathbf{A}_n), (\mathbf{A}_m')\) 是拓扑空间 \(\mathbf{X}\) 的子集的两个序列。假设对每对 \((\mathbf{A}_n, \mathbf{A}_m')\) 有一个波莱尔集 \(\mathbf{B}_{nm}\) （ \(\mathbf{X}\) 中的），使得 \(\mathbf{B}_{nm} \supset \mathbf{A}_n\) 且 \(\mathbf{B}_{nm} \cap \mathbf{A}_m' = \emptyset\) 。则存在一个波莱尔集 \(\mathbf{B}\) （ \(\mathbf{X}\) 中的），它包含 \(\bigcup_{n} \mathbf{A}_n\) 且不交 \(\bigcup_{m} \mathbf{A}_m'\) 。

因为波莱尔集 \(\mathbf{B} = \bigcup_{n}\left(\bigcap_{m}\mathrm{B}_{nm}\right)\) 满足这些条件。

引理 5. 设 X 是豪斯多夫空间，A 、A' 是 X 的两个不相交的苏斯林子空间。则存在 X 的一个波莱尔集 B ，使得 B ⊃ A 且 B ∩ A' = ∅。

由第 5 号的命题 13 ，存在两个筛 \(\mathbf{C}\) 、\(\mathbf{C}'\) 与连续映射 \(f\) （ \(\mathbf{L}(\mathbf{C})\) 到 \(\mathbf{A}\) 上的）、\(f'\) （ \(\mathbf{L}(\mathbf{C}')\) 到 \(\mathbf{A}'\) 上的），它们用第 5 号中描述的方法构造。对每个 \(n \geqslant 0\) 与每个 \(c \in \mathbf{C}_n\) ，令 \(q_n(c)\) 表示 \(\mathbf{L}(\mathbf{C})\) 中由一切序列 \((c_k)_{k \geqslant 0}\) （满足 \(c_n = c\) 者）所成的子空间；\(q_n(c)\) 是 \(\mathbf{L}(\mathbf{C})\) 的闭子空间。对每个 \(\gamma = (c_n) \in \mathbf{L}(\mathbf{C})\) ，闭集序列 \(q_n(c_n)\) 是递减的，且构成以 \(\gamma\) 为极限的滤基。此外，对每个 \(c \in \mathbf{C}_n\) ，诸集 \(q_{n+1}(d)\) （其中 \(d\) 跑遍集 \(\overline{p_n}(c)\) ，此集在 \(\mathbf{C}_{n+1}\) 中）构成 \(q_n(c)\) 的一个分划。对筛 \(c'\) 作类似的记号定义。

我们将用反证法，假设每个包含 A 的波莱尔集都与 \(A'\) 相交。首先，由引理 4 与筛分的定义，存在 \(c_{0} \in C_{0}\) 与 \(c_{0}' \in C_{0}'\) ，使得每个包含 \(f(q_{0}(c_{0}))\) 的波莱尔集都与 \(f'(q_{0}'(c_{0}'))\) 相交。然后我们可以对 n 用归纳法定义

\[
\gamma = (c _ {n}) \in \mathrm{L} (\mathrm{C}) \quad \text { and } \quad \gamma^ {\prime} = (c _ {n} ^ {\prime}) \in \mathrm{L} (\mathrm{C} ^ {\prime})
\]

如下：设 \(c_i\) 与 \(c_i'\) 已对 \(i < n\) 定义，使得对每个指标 \(i < n\) ，每个包含 \(f(q_i(c_i))\) 的波莱尔集都与 \(f'(q_i'(c_i'))\) 相交；把引理 4 与筛分的定义应用于集 \(f(q_{n-1}(c_{n-1}))\) 与 \(f'(q_{n-1}'(c_{n-1}'))\) ，即见存在 \(c_n \in \mathbf{C}_n\) 与 \(c_n' \in \mathbf{C}_n'\) ，使得 \(p_{n-1}(c_n) = c_{n-1}\) 且 \(p_{n-1}(c_n') = c_{n-1}'\) 。

并使得每个包含 \(f(q_n(c_n))\) 的波莱尔集都与 \(f'(q_n'(c_n'))\) 相交。而序列 \(f(q_n(c_n))\) 收敛于一点 \(a = f(\gamma) \in \mathbf{A}\) ，序列 \(f'(q_n'(c_n'))\) 收敛于一点 \(a' = f'(\gamma') \in \mathbf{A}'\) 。由于 \(\mathbf{A} \cap \mathbf{A}' = \emptyset\) 且 \(\mathbf{X}\) 是豪斯多夫的，存在一个闭邻域 \(\mathbf{V}\) （ \(a\) 的），它不含 \(a'\) ，从而对充分大的 \(n\) ，\(\mathbf{V}\) 包含 \(f(q_n(c_n))\) 且不交 \(f'(q_n'(c_n'))\) 。这是矛盾，因为 \(\mathbf{V}\) 是波莱尔集。

为证定理 2 ，令 \(\mathbf{Y}_n\) 表示诸集 \(\mathbf{X}_i\) （满足 \(i \neq n\) 者）的并；则 \(\mathbf{Y}_n\) 是 \(\mathbf{X}\) 的苏斯林子空间（第 2 号，命题 8）。由引理 5 ，对每个指标 \(n\) 存在波莱尔集 \(\mathbf{B}_n'\) ，它包含 \(\mathbf{X}_n\) 且不交 \(\mathbf{Y}_n\) 。令 \(\mathbf{B}_n\) 为 \(\mathbf{B}_n'\) 与 \(\bigcap_{i < n} (\mathbf{X} - \mathbf{B}_i')\) 的交。则诸 \(\mathbf{B}_n\) 是波莱尔集，两两不相交，且 \(\mathbf{B}_n \supset \mathbf{X}_n\) 对每个 \(n\) 成立。

推论. 若可度量化空间的一个可数分划由苏斯林集组成，则这些集是波莱尔集。特别地，可度量化空间中余集为苏斯林集的每个苏斯林集是波莱尔集。

